% \instlogo{} is deliberately left unset -- see tiet-question-paper.cls
% for the patch that makes that work correctly (upstream hard-errors
% otherwise). Call \instlogo{path/to/logo.pdf} to add your own crest.
\documentclass[11pt,a4paper,onecolumn]{tiet-question-paper}
\date{8 December 2026}
\institute{Riverside Institute of Technology}
\schoolordepartment{%
  Computer Science \& Engineering Department}
\examname{%
  End Semester Examination}
\coursecode{UCS505}
\coursename{Computer Graphics}
\timeduration{3 hours}
\maxmarks{45}
\faculty{ANG,AMK,HPS}

\begin{document}
\maketitle
\textbf{Instructions:}
\begin{enumerate}
\item Attempt any 5 questions;
\item Attempt all the subparts of a question at one
  place.
\end{enumerate}
\bvrhrule[0.4pt]
\begin{enumerate}
\item
  \begin{enumerate}
  \item Given the control polygon
    $\textbf{b}_0, \textbf{b}_1, \textbf{b}_2,
    \textbf{b}_3$ of a Cubic Bezier curve; determine
    the coordinates for parameter values
    $\forall t\in T$. \hfill [7 marks]

    \begin{align*}
      T \equiv
      & \{0, 0.15, 0.35, 0.5, 0.65, 0.85, 1\} \\
      \begin{bmatrix}
        \textbf{b}_0 &\textbf{b}_1& \textbf{b}_2& \textbf{b}_3
      \end{bmatrix} \equiv
      & \begin{bmatrix}
        1&2&4&3\\ 1&3&3&1
      \end{bmatrix}
    \end{align*}
  \item Explain the role of the convex hull property in
    bounding a Bezier curve. \hfill[2 marks]
  \end{enumerate}
\end{enumerate}
\bvrhrule[0.4pt]
\begin{enumerate}[resume]
\item
  \begin{enumerate}
  \item Describe the continuity conditions
    ($C^0$, $C^1$, $C^2$) for curvilinear geometry.
    \hfill[5 marks]
  \item Define formally, a B-Spline curve. \hfill [2
    marks]
  \item How is a Bezier curve different from a B-Spline
    curve? \hfill [3 marks]
  \end{enumerate}
\end{enumerate}
\bvrhrule[0.4pt]
\begin{enumerate}[resume]
\item
  \begin{enumerate}
  \item Given a triangle, with vertices defined by
    column vectors of $P$; find its vertices after
    reflection across the XZ plane. \hfill [3 marks]
    \begin{align*}
      P\equiv
      &\begin{bmatrix}
        3&6&5 \\ 4&4&6 \\ 1&2&3
      \end{bmatrix}
    \end{align*}
  \item Given a pyramid with vertices defined by the
    column vectors of $P$, and an axis of rotation $A$
    with direction $\textbf{v}$ and passing through
    $\textbf{p}$, find the coordinates of the vertices
    after rotation about $A$ by an angle of
    $\theta=\pi/4$.\hfill [6 marks]
    \begin{align*}
      P\equiv
      &\begin{bmatrix}
        0&1&0&0 \\ 0&0&1&0 \\0&0&0&1
      \end{bmatrix} \\
      \begin{bmatrix}
        \mathbf{v} & \mathbf{p}
      \end{bmatrix}\equiv
      &\begin{bmatrix}
        0&0 \\1&1\\1&0
      \end{bmatrix}
    \end{align*}
  \end{enumerate}
\end{enumerate}
\bvrhrule[0.4pt]
\begin{enumerate}[resume]
\item
  \begin{enumerate}
    \item Explain the two winding-number rules used for
      inside/outside tests in polygon filling. \hfill [4 marks]
    \item Explain the working principle of a
      CRT display. \hfill [5 marks]
  \end{enumerate}
\end{enumerate}
\bvrhrule[0.4pt]
\begin{enumerate}[resume]
\item
  \begin{enumerate}
  \item Given a projection plane $P$ defined by normal
    $\textbf{n}$ and a reference point $\textbf{a}$,
    and the centre of projection as $\mathbf{p}_0$,
    find the perspective projection of the point
    $\textbf{x}$ onto $P$. \hfill [5 marks]
    \begin{align*}
      \begin{bmatrix}
        \mathbf{a}&\mathbf{n}&\mathbf{p}_0&\mathbf{x}
      \end{bmatrix}\equiv
      &
        \begin{bmatrix}
          3&-1&1&8\\4&2&1&10\\5&-1&3&6
        \end{bmatrix}
    \end{align*}
  \item State the two main differences between a
    perspective projection and a parallel (orthographic)
    projection. \hfill [2 marks]
  \end{enumerate}
\end{enumerate}
\bvrhrule[0.4pt]
\end{document}
